\scanpage{326}\mspage{242}
% l. 1–3
d'après l'implication $b''\Rightarrow a$\inserted{$\Rightarrow c)$} de la \unsure{partie} $1^\circ)$ (cf page 235) et comme $b''_{\unsure{\pi|\kappa}}\Rightarrow b''$, on sait \deleted{\unsure{alors}} \unsure{déjà} que $f$ satisfait à $\mathrm{F}3$.
% l. 4–7
En l'appliquant à $j,k$ et $f(j)\longrightarrow f(k)$, OPS que $f(j)\longrightarrow f(k)$ provient de $v:j\unsure{\mathbin{-}|\longrightarrow}k$. On aura donc $\deleted{f(\unsure{vu_{\alpha}|vu'_{\alpha}})=f(}\; f(v)f(u_{\alpha})=f(v)f(u_{\beta})$ $\forall\alpha,\beta\in A$, donc $f(vu_{\alpha})=f(vu_{\beta})$.
% l. 8–11
Comme $f$ est fidèle, cela \unsure{indique|implique} que $vu_{\alpha}=vu_{\beta}$, i.e. $v$ égalise les $u_{\alpha}$. D'ailleurs, $v$ peut être considéré comme une flèche \uline{dans} $i'\backslash I$, de $(j,v')$ dans $(k,f(v)v')$. Qed.

% l. 12–15
\uline{4.9.6.10.} Quand $I$ est \uline{ordonné}, alors les quatre conditions $\unsure{a}_{\pi}\unsure{,|\Rightarrow}b_{\pi},c_{\pi},b''_{\pi}$ sont équivalentes, car les $i'\backslash I$ sont \unsure{ordonnées|ordonnés}, et il est clair qu'\uline{alors} $b''_{\pi}\unsure{\Longleftrightarrow|\Longrightarrow}\unsure{b^{\mathrm{s}}_{\pi}|b_{\pi}}$\unsure{.}

% l. 16
\deleted{\uline{4.9.7} Démonstration de 4.9.\unsure{3|2}}
% archive sheet number: 325
